% GENERATED FILE. Edit canonical lessons, then run npm run generate:books.
% canonical-source-hash: fnv1a64:f799345af1afd039
% canonical-lessons: SA-C125-nakulah, SA-C125-hrasvah, SA-C125-mrduh, SA-C125-suddhah, SA-C125-malinah

\chapter{Mongoose, Short, Soft, Clean, Dirty}
\label{ch:sa-mongoose-short-soft-clean-dirty}
\hlchaptermodality{125}

\begin{chapteropening}
\textbf{By the end of this chapter:} \emph{I can say a mongoose, short, soft, clean and dirty.}

Complete the last lesson of chapter 125: I can say a mongoose, short, soft, clean and dirty.
\end{chapteropening}

\section[\texorpdfstring{\sk{नकुलः} (nakulaḥ)}{नकुलः (nakulaḥ)}]{\sk{नकुलः} (nakulaḥ) — a mongoose}

\label{lesson:SA-C125-nakulah}



\begin{quote}
Before the new one: say the Sanskrit for a cuckoo, then the Sanskrit for a sparrow.
\end{quote}

\subsection*{You'll want to know: \sk{नकुलः}}
\textbf{\sk{नकुलः}} — \emph{nakulaḥ} — \textquotedblleft{}a mongoose\textquotedblright{}.

\begin{grammarlens}[title={What you've built}]
One more word for this chapter.
\end{grammarlens}

\subsection*{Your turn}
\begin{itemize}
\raggedright
  \item \emph{Say it:} \emph{nakulaḥ}
  \item \emph{Say it:} \emph{nakulaḥ}, once more
  \item \emph{Say it:} say \emph{nakulaḥ}
  \item \emph{Recall:} say the Sanskrit for a cuckoo, then the Sanskrit for a sparrow, then say \emph{nakulaḥ} again
\end{itemize}

\begin{tcolorbox}[breakable,colback=teal!4,colframe=teal!35!black,title={Before you move on}]
What does \sk{नकुलः} mean? (\textquotedblleft{}A mongoose\textquotedblright{}.) Say it once more.
\end{tcolorbox}
\section[\texorpdfstring{\sk{ह्रस्वः} (hrasvaḥ)}{ह्रस्वः (hrasvaḥ)}]{\sk{ह्रस्वः} (hrasvaḥ) — short}

\label{lesson:SA-C125-hrasvah}



\begin{quote}
Before the new one: say the Sanskrit for a sparrow, then the Sanskrit for a mongoose.
\end{quote}

\subsection*{You'll want to know: \sk{ह्रस्वः}}
\textbf{\sk{ह्रस्वः}} — \emph{hrasvaḥ} — \textquotedblleft{}short\textquotedblright{}.

\begin{grammarlens}[title={What you've built}]
One more word for this chapter.
\end{grammarlens}

\subsection*{Your turn}
\begin{itemize}
\raggedright
  \item \emph{Say it:} \emph{hrasvaḥ}
  \item \emph{Say it:} \emph{hrasvaḥ}, once more
  \item \emph{Say it:} say \emph{hrasvaḥ}
  \item \emph{Recall:} say the Sanskrit for a sparrow, then the Sanskrit for a mongoose, then say \emph{hrasvaḥ} again
\end{itemize}

\begin{tcolorbox}[breakable,colback=teal!4,colframe=teal!35!black,title={Before you move on}]
What does \sk{ह्रस्वः} mean? (\textquotedblleft{}Short\textquotedblright{}.) Say it once more.
\end{tcolorbox}
\section[\texorpdfstring{\sk{मृदुः} (mṛduḥ)}{मृदुः (mṛduḥ)}]{\sk{मृदुः} (mṛduḥ) — soft}

\label{lesson:SA-C125-mrduh}



\begin{quote}
Before the new one: say the Sanskrit for a mongoose, then the Sanskrit for short.
\end{quote}

\subsection*{You'll want to know: \sk{मृदुः}}
\textbf{\sk{मृदुः}} — \emph{mṛduḥ} — \textquotedblleft{}soft\textquotedblright{}.

\begin{grammarlens}[title={What you've built}]
One more word for this chapter.
\end{grammarlens}

\subsection*{Your turn}
\begin{itemize}
\raggedright
  \item \emph{Say it:} \emph{mṛduḥ}
  \item \emph{Say it:} \emph{mṛduḥ}, once more
  \item \emph{Say it:} say \emph{mṛduḥ}
  \item \emph{Recall:} say the Sanskrit for a mongoose, then the Sanskrit for short, then say \emph{mṛduḥ} again
\end{itemize}

\begin{tcolorbox}[breakable,colback=teal!4,colframe=teal!35!black,title={Before you move on}]
What does \sk{मृदुः} mean? (\textquotedblleft{}Soft\textquotedblright{}.) Say it once more.
\end{tcolorbox}
\section[\texorpdfstring{\sk{शुद्धः} (śuddhaḥ)}{शुद्धः (śuddhaḥ)}]{\sk{शुद्धः} (śuddhaḥ) — clean, pure}

\label{lesson:SA-C125-suddhah}



\begin{quote}
Before the new one: say the Sanskrit for short, then the Sanskrit for soft.
\end{quote}

\subsection*{You'll want to know: \sk{शुद्धः}}
\textbf{\sk{शुद्धः}} — \emph{śuddhaḥ} — \textquotedblleft{}clean, pure\textquotedblright{}.

\begin{grammarlens}[title={What you've built}]
One more word for this chapter.
\end{grammarlens}

\subsection*{Your turn}
\begin{itemize}
\raggedright
  \item \emph{Say it:} \emph{śuddhaḥ}
  \item \emph{Say it:} \emph{śuddhaḥ}, once more
  \item \emph{Say it:} say \emph{śuddhaḥ}
  \item \emph{Recall:} say the Sanskrit for short, then the Sanskrit for soft, then say \emph{śuddhaḥ} again
\end{itemize}

\begin{tcolorbox}[breakable,colback=teal!4,colframe=teal!35!black,title={Before you move on}]
What does \sk{शुद्धः} mean? (\textquotedblleft{}Clean, pure\textquotedblright{}.) Say it once more.
\end{tcolorbox}
\section[\texorpdfstring{\sk{मलिनः} (malinaḥ)}{मलिनः (malinaḥ)}]{\sk{मलिनः} (malinaḥ) — dirty}

\label{lesson:SA-C125-malinah}



\begin{quote}
Before the new one: say the Sanskrit for soft, then the Sanskrit for clean.
\end{quote}

\subsection*{You'll want to know: \sk{मलिनः}}
\textbf{\sk{मलिनः}} — \emph{malinaḥ} — \textquotedblleft{}dirty\textquotedblright{}.

\begin{grammarlens}[title={What you've built}]
One more word for this chapter.
\end{grammarlens}

\subsection*{Your turn}
\begin{itemize}
\raggedright
  \item \emph{Say it:} \emph{malinaḥ}
  \item \emph{Say it:} \emph{malinaḥ}, once more
  \item \emph{Say it:} say \emph{malinaḥ}
  \item \emph{Recall:} say the Sanskrit for soft, then the Sanskrit for clean, then say \emph{malinaḥ} again
\end{itemize}

\begin{tcolorbox}[breakable,colback=teal!4,colframe=teal!35!black,title={Before you move on}]
What does \sk{मलिनः} mean? (\textquotedblleft{}Dirty\textquotedblright{}.) Say it once more.
\end{tcolorbox}
